\documentclass[10pt,letterpaper]{article}
\usepackage{neurips_2024}

% Extra packages used in this paper
\usepackage{xspace}
\usepackage{url}
\usepackage{amsmath}
\usepackage{booktabs}
\usepackage{tabularx}
\usepackage{multirow}
\usepackage{makecell}
\usepackage{enumitem}
\usepackage{listings}
\usepackage{parskip}
\usepackage{float}

\setlength{\parskip}{4pt}
\setlength{\parindent}{12pt}

% Convenience macros
\newcommand{\gridcast}{\textsc{GridCast}\xspace}
\newcommand{\graphcast}{\textsc{GraphCast}\xspace}
\newcommand{\gencast}{\textsc{GenCast}\xspace}
\newcommand{\eg}{\textit{e.g.,}\xspace}
\newcommand{\ie}{\textit{i.e.,}\xspace}

% ============================================================
\title{GridCast: Probabilistic Grid Stress Forecasting\\
       for Dynamic Data Center Power Allocation}

\subtitle{Project Proposal $\vert$ Draft for Review}

\author{%
  \textit{[Author names omitted for review]}
}

\date{}

% ============================================================
\begin{document}

\maketitle

% ============================================================
\begin{abstract}
Data centers increasingly strain power grid infrastructure, yet flexible power contracts remain
underutilised because utilities cannot forecast grid stress with enough confidence to underwrite
contractual commitments. \gridcast is a probabilistic deep-learning pipeline that combines
\graphcast and \gencast (Google DeepMind's deterministic and ensemble weather models) with
regional grid context from the EIA API and node-level stress signals from Grid~Status
locational marginal price (LMP) data. Together, these three layers produce
confidence-bounded grid stress estimates at the transmission-node level, giving utilities the
uncertainty-aware signal needed to offer dynamic power contracts. A key advantage is
\graphcast's trainability: unlike physics-locked NWP models, it can be fine-tuned on recent
observational data to adapt to shifting climate trends.
\end{abstract}

% ============================================================
\section{Introduction and Problem Statement}

The power grid is overbuilt by design, and most excess capacity sits idle nearly every day.
Under flexible power contracts, data centers would commit to drawing less during grid stress
events---compensating through workload shifting or graceful degradation---in exchange for
access to that idle capacity at preferential rates.

The bottleneck is not a lack of forecasting. Utilities already use meteorological inputs for
demand forecasting. The problem is the \emph{absence of forecasting confident enough to
underwrite contractual commitments}. Grid stress emerges from compounding tail-risk
scenarios: a heatwave extending longer than forecast, a cold snap of unexpected severity,
or high demand coinciding with renewable underperformance. Deterministic point forecasts
cannot communicate the probability of such scenarios. Without uncertainty bounds, no utility
can confidently offer a tiered contract, and no data center can reliably plan around one.

\gridcast addresses this by producing confidence-bounded, node-level stress estimates over
a 10-day horizon.

% ============================================================
\section{Why AI Weather Forecasting}

\graphcast operates at $0.25$-degree global resolution and can be queried over precise
coordinates, enabling node-level reasoning rather than regional averaging. As a learned model
trained on the ERA5 reanalysis archive, it can be fine-tuned on recent data to adapt to
shifting climate patterns, which physics-locked NWP systems cannot do. It also produces
10-day global forecasts in under one minute, enabling near-real-time updates for hourly
contract triggers.

\gencast extends \graphcast by generating an \emph{ensemble} of plausible futures rather
than a single deterministic trajectory. A utility does not need to know a heatwave \emph{will}
occur; it needs to know the probability and severity range. This ensemble output is what makes
the downstream stress signal actionable for contract design.

% ============================================================
\section{Proposed Approach}

\subsection{Three-Layer Pipeline}

\paragraph{Layer~1: Probabilistic Weather Forecasting (\graphcast + \gencast).}
\graphcast forecasts key atmospheric variables at 6-hour intervals over a 10-day horizon.
ERA5 reanalysis serves as the training ground truth; at inference time, \graphcast forecasts
replace ERA5. \gencast wraps \graphcast's output with a diffusion-based ensemble,
producing a distribution over plausible futures with explicit confidence intervals that
propagate through the downstream layers.

\paragraph{Layer~2: Regional Grid Context (EIA API).}
EIA hourly demand and fuel-mix data from Form EIA-930 (\texttt{/v2/electricity/rto/}) are
used to train the weather-to-load relationship for each Balancing Authority. EIA data also
captures supply-side context: renewable penetration, available headroom, and the
forecast-versus-actual demand delta. These variables are essential for identifying compounding
scenarios where demand spikes coincide with renewable underperformance.

\paragraph{Layer~3: Node-Level Stress Signal (Grid Status LMP).}
Locational Marginal Price at a target node is the wholesale market's real-time expression of
congestion and supply tightness at that specific location, providing the node-level resolution
that EIA cannot supply. Historical nodal LMP from the open-source \texttt{gridstatus} Python
library is used to train a mapping from regional stress to node-level price behaviour; day-ahead
LMP serves as a forward signal at inference time.

% ------------------------------------------------------------
\subsection{Key Input Variables}

Table~\ref{tab:inputs} summarises the input features, their sources, and their roles in the
pipeline.

\begin{table}[H]
\centering
\caption{Key input variables, data sources, and pipeline roles.}
\label{tab:inputs}
\renewcommand{\arraystretch}{1.25}
\begin{tabularx}{\linewidth}{@{}lllX@{}}
\toprule
\textbf{Category} & \textbf{Variable} & \textbf{Source} & \textbf{Role} \\
\midrule
Weather & 2\,m Temperature & \graphcast / ERA5 & Primary demand driver \\
Weather & 10\,m Wind Speed & \graphcast / ERA5 & Renewable generation proxy \\
Weather & Ensemble Spread & \gencast & Uncertainty quantification \\
\midrule
Grid (EIA) & Hourly Actual Demand & EIA \texttt{/rto/region-data} & Training target variable \\
Grid (EIA) & Generation by Fuel Type & EIA \texttt{/rto/fuel-type} & Renewable underperformance \\
\midrule
Grid (LMP) & Real-Time Nodal LMP & Grid Status & Node-level stress signal \\
Grid (LMP) & LMP Congestion Component & Grid Status & Transmission constraint \\
\midrule
Engineered & Hour / Day / Month & Derived & Demand periodicity \\
Engineered & Demand / Capacity Ratio & Derived from EIA & Interpretable stress metric \\
\bottomrule
\end{tabularx}
\end{table}

% ------------------------------------------------------------
\subsection{Training and Output}

The weather-to-grid-stress mapping is trained on ERA5 paired with EIA historical data from
2018 to present, the period over which both are available at compatible resolution. Training on
reanalysis rather than forecast outputs is essential: forecasts carry errors that would otherwise
propagate into the learned mapping.

At inference time, \graphcast and \gencast replace ERA5, and the pipeline produces a
confidence-bounded probability distribution over grid stress at a target node over the 10-day
horizon. This is operationalised as a \emph{recommended power allocation percentage}: the
share of excess capacity at a given node that can safely be committed to a data center
operator, conditioned on forecast confidence.

% ============================================================
\section{Validation}

The pipeline is validated against historical EIA load data and nodal LMP during documented
stress events. Three well-documented cases are used, covering distinct failure modes and
grid topologies:

\begin{itemize}
  \item \textbf{2021 Texas Winter Storm.} Cold-driven demand surge; extreme exceedance of
        forecast demand in ERCOT territory.
  \item \textbf{2022 Pacific Northwest Heat Dome.} Sustained multi-day heat event testing
        the model's ability to maintain calibration over a prolonged horizon.
  \item \textbf{2023 PJM Summer Stress Events.} Compounding of high load with renewable
        underperformance across a large, highly interconnected grid topology.
\end{itemize}

% ============================================================
\section{Value Proposition}

This project demonstrates that a probabilistic, node-level forecasting layer for dynamic power
contracts is technically viable today using existing open-source tools and public data.

\begin{itemize}
  \item \textbf{Utility Providers.} Node-level probabilistic forecasts let utilities offer tiered
        contracts with confidence, converting idle excess capacity into revenue.
  \item \textbf{Grid Operators (PJM, CAISO, MISO).} Better advance forecasting reduces
        costly last-minute interventions and improves visibility into renewable supply
        fluctuations.
  \item \textbf{Data Centers.} Forecast-driven contracts unlock excess grid capacity, reducing
        operational costs and enabling more sustainable scaling of AI compute.
  \item \textbf{Hardware Providers (\eg NVIDIA).} Power availability is a primary bottleneck
        to data center expansion. Unlocking excess grid capacity has clear strategic value for
        businesses whose growth depends on compute scaling.
\end{itemize}

% ============================================================
\appendix
\section{Data Sources}
\label{app:data}

\subsection{EIA API v2: Hourly Electric Grid Monitor}

\textbf{Source:} U.S.\ Energy Information Administration, Form EIA-930. Free public API at
\url{api.eia.gov/v2/electricity/rto/} (free key registration at \url{eia.gov/opendata}). Data
are organised by Balancing Authority at hourly resolution, with historical depth from July 2015
for demand and generation totals and July 2018 for generation by fuel type.

\begin{lstlisting}[language=bash, caption={Example EIA API query for PJM hourly demand.}, label={lst:eia}]
GET https://api.eia.gov/v2/electricity/rto/region-data/data/
  ?frequency=hourly
  &data[]=value
  &facets[respondent][]=PJM
  &facets[type][]=D
  &sort[0][column]=period
  &api_key=YOUR_KEY
\end{lstlisting}

\begin{table}[H]
\centering
\caption{EIA \texttt{/rto/region-data} sample --- PJM, 27 July 2023.}
\label{tab:eia-demand}
\renewcommand{\arraystretch}{1.2}
\begin{tabular}{@{}llllrr@{}}
\toprule
\textbf{Period} & \textbf{Resp.} & \textbf{Type} & \textbf{Type Name} & \textbf{Value} & \textbf{Units} \\
\midrule
2023-07-27T17 & PJM & D  & Demand           & 152,340 & MWh \\
2023-07-27T17 & PJM & DF & Forecast Demand  & 148,900 & MWh \\
2023-07-27T17 & PJM & NG & Net Generation   & 149,820 & MWh \\
2023-07-27T17 & PJM & TI & Total Interchange& $-$2,520 & MWh \\
\bottomrule
\end{tabular}
\end{table}

The forecast-versus-actual delta (152,340 vs.\ 148,900\,MWh) on a peak summer afternoon
illustrates the surprise stress deviation the pipeline is designed to anticipate \emph{before}
it occurs.

\begin{table}[H]
\centering
\caption{EIA \texttt{/rto/fuel-type-data} sample --- CAISO, 27 July 2023.}
\label{tab:eia-fuel}
\renewcommand{\arraystretch}{1.2}
\begin{tabular}{@{}lllllrr@{}}
\toprule
\textbf{Period} & \textbf{Resp.} & \textbf{Fuel} & \textbf{Fuel Name} & \textbf{Value} & \textbf{Units} \\
\midrule
2023-07-27T15 & CISO & SUN & Solar       &  9,842 & MWh \\
2023-07-27T15 & CISO & NG  & Natural Gas & 14,560 & MWh \\
2023-07-27T18 & CISO & SUN & Solar       &  3,410 & MWh \\
2023-07-27T18 & CISO & NG  & Natural Gas & 18,940 & MWh \\
\bottomrule
\end{tabular}
\end{table}

The drop in solar and corresponding rise in gas dispatch between 15:00 and 18:00 illustrates
the afternoon duck curve---a daily structural stress window in CAISO territory.

% ------------------------------------------------------------
\subsection{Grid Status: Nodal LMP Data}

\textbf{Source:} \url{gridstatus.io}. Open-source Python library (\texttt{pip install
gridstatus}) for recent data; commercial API tier for full historical depth. Covers PJM, CAISO,
MISO, ERCOT, ISO-NE, NYISO, and SPP at 5-minute real-time and hourly day-ahead resolution.
LMP decomposes into energy, congestion, and loss components. The congestion component
isolates transmission bottlenecks and is the primary node-level stress signal used in
\gridcast.

\begin{lstlisting}[language=Python, caption={Fetching PJM real-time LMP via \texttt{gridstatus}.}, label={lst:gridstatus}]
import gridstatus

pjm = gridstatus.PJM()
df  = pjm.get_lmp(
    date="2023-07-27",
    market="REAL_TIME_5_MIN",
    locations="ALL"
)
\end{lstlisting}

\begin{table}[H]
\centering
\caption{Grid Status real-time LMP --- PJM, 27 July 2023 (peak hour). All prices in \$/MWh.}
\label{tab:lmp}
\renewcommand{\arraystretch}{1.2}
\begin{tabular}{@{}llrrrr@{}}
\toprule
\textbf{Time} & \textbf{Location} & \textbf{LMP} & \textbf{Energy} & \textbf{Congestion} & \textbf{Loss} \\
\midrule
2023-07-27 17:00 & Dominion Hub      & 142.30 & 98.40 & \textbf{38.70} & 5.20 \\
2023-07-27 17:00 & AEP-Dayton Hub    & 101.15 & 98.40 & 1.40  & 1.35 \\
2023-07-27 17:00 & Western Hub       &  99.80 & 98.40 & 0.20  & 1.20 \\
2023-07-27 17:00 & N.\ Illinois Hub  & 103.60 & 98.40 & 3.70  & 1.50 \\
\bottomrule
\end{tabular}
\end{table}

Dominion Hub (Northern Virginia, a major data center corridor) shows a congestion component
of \$38.70/MWh---nearly 40\% of total LMP---while other hubs show near-zero congestion at the
same system-wide energy price. This spatial divergence is the node-level signal \gridcast is
designed to forecast before the stress event occurs.

% ------------------------------------------------------------
\subsection{ERA5 and Data Integration}
\label{app:era5}

ERA5 (ECMWF global reanalysis, available at \url{cds.climate.copernicus.eu} via the
\texttt{cdsapi} Python library) provides the ground-truth weather input for training at
$0.25$-degree hourly resolution from 1940 to near-present. All three sources are aligned to
UTC before training. Grid Status LMP timestamps vary by ISO and must be converted.
Geographic alignment maps EIA Balancing Authority regions to ERA5 grid cells, and Grid
Status node coordinates (annotated for approximately 20,000 pricing points) to the nearest
ERA5 grid point.

\end{document}